\section{Boundary information after residual-time integration}
\label{sec:v17-boundary-information}

Residual-time integration preserves the smooth relative law of
Theorem~\ref{thm:v8-main-relative}, but it changes the singularity of the
finite-versus-boundary experiment.  Before integration the tangent success
density is positive up to a moving boundary.  After integration it vanishes
there to first order.  The resulting logarithmic information is a boundary
quantity.  The theorem below isolates that quantity before the billiard
specialization.

The one-dimensional change of rate caused by parameter-dependent support is
classical.  Akahira~\cite{Akahira1975a,Akahira1975b} and
Ibragimov--Has'minskii~\cite{IbragimovHasminskii} belong to the general
nonregular background, and Smith~\cite{Smith1985} explicitly treats the
linearly vanishing endpoint case.  Limits of experiments for models with
parameter-dependent support are also developed by
Hirano--Porter~\cite{HiranoPorter}.  We do not use these references as a
black-box proof: the form needed here has a moving regular hypersurface, an
intrinsic coefficient, and an independent failure atom, and we prove it
directly.

\subsection{A linearly vanishing density}
Let $U\subset\mathbb R^m$ be bounded and open.  For $|t|<t_0$ let $a_t$
and $w_t$ be real functions which are $C^3$ jointly on a neighborhood of
$[-t_0,t_0]\times\overline U$.  Assume that $a_t$ is bounded above and
bounded away from zero, that $w_t$ is uniformly negative near
$\partial U$, and that
\[
 D=\{x\in U:w_0(x)>0\}\ne\varnothing,\qquad
 \Sigma=\{x\in U:w_0(x)=0\},\qquad
 |\nabla w_0|>0\quad\hbox{on }\Sigma.
\]
Thus $\overline D\subset U$ and $\Sigma=\partial D$ is a compact regular
hypersurface.  Decreasing $t_0$ if necessary, suppose
\begin{equation}\label{eq:v17-linear-density}
 f_t(x)=a_t(x)(w_t(x))_+,
 \qquad \int_U f_t(x)\,dx=1.
\end{equation}
Normalizing an initially unnormalized density is allowed: a common tubular
chart for the moving boundary makes its positive normalizing integral
$C^3$.  Put
\begin{equation}\label{eq:v17-boundary-form}
 v(x)=\left.\partial_t w_t(x)\right|_{t=0},\qquad
 \mathcal I_\Sigma=
 \int_\Sigma\frac{a_0(x)v(x)^2}{|\nabla w_0(x)|}\,d\sigma(x).
\end{equation}
We use
\[
 H^2(P,Q)=\int(\sqrt{dP/d\nu}-\sqrt{dQ/d\nu})^2\,d\nu,
 \qquad \|P-Q\|_{\rm TV}=\sup_B|P(B)-Q(B)|,
\]
so $0\le H^2\le2$, and write $\Phi$ for the standard normal distribution
function.

\begin{theorem}[Logarithmic boundary information]
\label{thm:v17-boundary-information}
Under \eqref{eq:v17-linear-density},
\begin{equation}\label{eq:v17-hellinger}
 H^2(f_t\,dx,f_0\,dx)
 =\frac{\mathcal I_\Sigma}{4}t^2\log\frac1{|t|}+O(t^2),
 \qquad t\to0.
\end{equation}
Assume $\mathcal I_\Sigma>0$.  Let $t_n\to0$, $t_n\ne0$, let
$0<p_n\le1$, and consider $n$ independent observations from
\begin{equation}\label{eq:v17-thinning}
 P_{n,0}=(1-p_n)\delta_\dagger+p_nf_0\,dx,
 \qquad
 P_{n,1}=(1-p_n)\delta_\dagger+p_nf_{t_n}\,dx,
\end{equation}
where the failure symbol $\dagger$ is retained.  If
\begin{equation}\label{eq:v17-information-scale}
 np_nt_n^2\log(1/|t_n|)\longrightarrow b\in[0,\infty],
\end{equation}
then
\begin{equation}\label{eq:v17-general-testing}
 \|P_{n,1}^{\otimes n}-P_{n,0}^{\otimes n}\|_{\rm TV}
 \longrightarrow 2\Phi(\sqrt{\mathcal I_\Sigma b}/2)-1,
\end{equation}
with the right-hand side interpreted as one when $b=\infty$.

If $0<b<\infty$, the log likelihood of the absolutely continuous part of
$P_{n,1}^{\otimes n}$ relative to $P_{n,0}^{\otimes n}$ converges under
the latter law to
\begin{equation}\label{eq:v17-general-likelihood}
 \mathcal N(-\mathcal I_\Sigma b/2,\mathcal I_\Sigma b),
\end{equation}
and the minimum equal-prior testing error tends to
$\Phi(-\sqrt{\mathcal I_\Sigma b}/2)$.

If $b<\infty$ (including $b=0$), the probability that at least one of the
$n$ observations falls in a support-exclusive region tends to zero under
either hypothesis.  No support-exclusivity assertion is made under the
sole assumption $b=\infty$; the conclusion
\eqref{eq:v17-general-testing} in that regime follows from Hellinger
affinity and remains valid even when $np_nt_n^2$ does not tend to zero.
\end{theorem}

The last paragraph is a scope distinction, not a weakening of the
supercritical theorem.  At finite information scale,
$np_nt_n^2=(np_nt_n^2\log(1/|t_n|))/\log(1/|t_n|)\to0$, so exclusive
support is negligible.  At supercritical scale that implication is false
in general, but it is unnecessary for total-variation separation.

\subsection{Coarea estimates}
For $x\in D$ define
\[
 S(x)=\frac{\dot a_0(x)}{a_0(x)}+\frac{v(x)}{w_0(x)}.
\]
It is integrable under $f_0\,dx$ but, when $\mathcal I_\Sigma>0$, has
infinite second moment.  For $h>0$ put
$D_h=\{w_0\ge h\}$ and $S_h=S\mathbf1_{D_h}$.

\begin{lemma}[Truncated score moments]\label{lem:v17-coarea}
As $h\downarrow0$,
\begin{align}
 \Pr_0\{0<w_0<h\}&=O(h^2),& \mathbb E_0S_h&=O(h),
 \label{eq:v17-strip}\\
 \mathbb E_0S_h^2&=\mathcal I_\Sigma\log(1/h)+O(1),
 \label{eq:v17-score-two}\\
 \mathbb E_0|S_h|^3&=O(h^{-1}),&
 \mathbb E_0|S_h|^4&=O(h^{-2}).
 \label{eq:v17-score-higher}
\end{align}
For $|t|$ sufficiently small, the $f_t$-mass of its support outside $D$
and the $f_0$-mass of the region where $f_t=0$ are both $O(t^2)$.
The estimates are uniform on compact families with the stated strict
positivity and regularity bounds.
\end{lemma}

\begin{proof}
Choose $h_0>0$ so that $|\nabla w_0|$ is bounded away from zero on
$|w_0|\le2h_0$.  The flow of
$\nabla w_0/|\nabla w_0|^2$ identifies the level sets $w_0=z$ for
$|z|\le h_0$ with $\Sigma$, with uniformly bounded $C^2$ Jacobians.
By coarea,
\[
 J(z)=\int_{\{w_0=z\}}\frac{a_0v^2}{|\nabla w_0|}\,d\sigma
      =\mathcal I_\Sigma+O(|z|).
\]
The collar contribution of $v^2/w_0^2$ to the second moment is therefore
$\int_h^{h_0}J(z)\,dz/z$.  The remaining terms in $S^2$ are integrable,
which proves \eqref{eq:v17-score-two}.  Since $f_0=O(w_0)$ in the collar,
\[
 \int_{\{0<w_0<h\}}f_0\,dx=O(h^2),\qquad
 \int_{D_h}|S|^r f_0\,dx
 \le C\left(1+\int_h^{h_0}z^{1-r}\,dz\right),\quad r=3,4.
\]
Differentiating $\int f_t=1$ at zero gives
$\int_D(\dot a_0w_0+a_0v)\,dx=0$, hence $\mathbb E_0S=0$.
The omitted collar integral of $Sf_0$ is $O(h)$, proving
\eqref{eq:v17-strip}.  Finally $|w_t-w_0|\le C|t|$, so both
support-exclusive regions lie in $|w_0|\le C|t|$; there the relevant
density is $O(|t|)$ and the collar volume is $O(|t|)$.  Their masses are
$O(t^2)$.
\end{proof}

\subsection{Proof of the information theorem}
\begin{proof}[Proof of Theorem~\ref{thm:v17-boundary-information}]
Write $t\ne0$ and $\ell=\log(1/|t|)$.  On $D_h$, whenever
$h\ge C_0|t|$ for a sufficiently large fixed $C_0$, Taylor expansion of
$a_tw_t$ before division by $a_0w_0$ gives
\begin{equation}\label{eq:v17-density-ratio}
 \frac{f_t}{f_0}=1+tS+t^2R_t,
 \qquad |R_t(x)|\le C(1+w_0(x)^{-1}).
\end{equation}
This expansion never differentiates a support indicator.

Take first $h=C_0|t|$.  The integral of $f_0+f_t$ outside the common part
$D_h$ is $O(t^2)$.  On $D_h$ set $z=tS+t^2R_t$.  With $C_0$ large and
then $t$ small, $|z|\le1/2$, and
$(\sqrt{1+z}-1)^2=z^2/4+O(|z|^3)$.  The coarea bounds give
\[
 \int_{D_h}z^2f_0\,dx
 =t^2\mathbb E_0S_h^2+O(|t|^3\log(1/|t|)),
 \qquad
 \int_{D_h}|z|^3f_0\,dx=O(t^2).
\]
Together with \eqref{eq:v17-score-two} this proves
\eqref{eq:v17-hellinger}.

Assume now $0<b<\infty$.  In the $n$th row write
\[
 t=t_n,\quad p=p_n,\quad \ell=\log(1/|t|),\quad
 L=\ell^{1/4},\quad h=|t|L.
\]
Represent a null observation by a Bernoulli success indicator $Y_i$ of
mean $p$ and, on success, an independent $X_i$ of density $f_0$.  With
probability tending to one all successful $X_i$ lie in $D_h$, because
\[
 Cnph^2=C(npt^2\ell)L^2/\ell=o(1).
\]
The same is true under the alternative.  In particular the total
support-exclusive probability is $O(npt^2)=o(1)$ in this finite critical
regime.

On the high-probability common-support event,
\begin{equation}\label{eq:v17-log-expansion}
 \log\frac{dP_{n,1}^{\otimes n}}{dP_{n,0}^{\otimes n}}
 =t\sum_{i=1}^nY_iS_h(X_i)
 -\frac{t^2}{2}\sum_{i=1}^nY_iS_h(X_i)^2+o_{P_{n,0}^{\otimes n}}(1).
\end{equation}
Indeed the summed $t^2R_t$ contribution is $O(npt^2)=o(1)$,
the cubic logarithmic remainder has expectation
$O(np|t|^3/h)=O((npt^2\ell)/(L\ell))=o(1)$, and replacing the quadratic
term by $t^2S_h^2$ produces $o(1)$ by the same coarea estimates.

The mean of the first sum in \eqref{eq:v17-log-expansion} is
$O(np|t|h)=o(1)$ and its variance tends to
$\mathcal I_\Sigma b$.  Each centered summand after multiplication by $t$
is $O(1/L)$, so Lindeberg is automatic.  The quadratic sum has mean
$\mathcal I_\Sigma b+o(1)$ and variance bounded by
\[
 Cnpt^4\mathbb E_0|S_h|^4
 \le Cnpt^4h^{-2}=C(npt^2\ell)/(L^2\ell)=o(1).
\]
This proves \eqref{eq:v17-general-likelihood}.

Let $Z_n$ be the log likelihood, with $Z_n=-\infty$ on the null-exclusive
set.  Even in the presence of an alternative singular part,
\[
 1-\|P_{n,1}^{\otimes n}-P_{n,0}^{\otimes n}\|_{\rm TV}
 =\mathbb E_{P_{n,0}^{\otimes n}}\min\{1,e^{Z_n}\}.
\]
The integrand is bounded and continuous on the extended real line.  If
$v=\mathcal I_\Sigma b$ and $G\sim\mathcal N(0,1)$, then
\[
 \mathbb E\min\{1,e^{\sqrt vG-v/2}\}=2\Phi(-\sqrt v/2),
\]
which proves the finite positive case of
\eqref{eq:v17-general-testing} and the testing-risk assertion.

For $b=0$ and $b=\infty$ use affinity instead.  The one-observation
affinity of the thinned pair is exactly
$1-pH^2(f_t,f_0)/2$.  Affinities multiply, and if $\mathcal A(P,Q)$ is
the affinity then
\[
 1-\mathcal A(P,Q)\le\|P-Q\|_{\rm TV}
 \le\sqrt{1-\mathcal A(P,Q)^2}.
\]
Equation~\eqref{eq:v17-hellinger} therefore gives TV $\to0$ when $b=0$
and TV $\to1$ when $b=\infty$, without any assumption on
$npt^2$.  If merely $b<\infty$, the support-exclusive probability is
$O(npt^2)=o(1)$ by the same division by $\ell$ used above.  This proves all
remaining statements.
\end{proof}

\begin{remark}[Intrinsic coefficient]\label{rem:v17-invariance}
Replace $w_t$ by $\widetilde w_t=c_tw_t$, $c_t>0$ smooth, and $a_t$ by
$a_t/c_t$.  On $\Sigma$,
$\partial_t\widetilde w_t|_0=c_0v$ and
$|\nabla\widetilde w_0|=c_0|\nabla w_0|$, so
\eqref{eq:v17-boundary-form} is unchanged.  Equivalently, if
$V_\nu=-v/|\nabla w_0|$ is the normal velocity of the support boundary,
\[
 \mathcal I_\Sigma=
 \int_\Sigma a_0|\nabla w_0|V_\nu^2\,d\sigma.
\]
If this integral vanishes, the theorem supplies only $H^2=O(t^2)$; no
nontrivial logarithmic critical scale is asserted.  When it is positive,
the smooth interior Fisher contribution is lower order.
\end{remark}
